%% 3.1.9 Resumo em língua portuguesa
%% Elemento obrigatório (Figura 14).
%% Consiste na apresentação sucinta dos pontos relevantes do texto,
%% em um único parágrafo. O resumo deve conter entre 150 e 500 pala-
%% vras e fornecer uma visão rápida e clara dos objetivos, da metodologia,
%% dos resultados e das conclusões do trabalho. Na elaboração do resumo,
%% deve-se usar o verbo na voz ativa, na terceira pessoa do singular.

%% Fonte -> TNR ou Arial, corpo 12.
%% A palavra RESUMO deve aparecer em letras maiúsculas
%% e em negrito.
%% O uso de itálico é permitido em palavras estrangeiras.
%% O uso de letras maiúsculas nas palavras-chave
%% restringe-se ao início da palavra, em nomes próprios
%% e siglas, se for o caso.

%% Alinhamento -> A palavra RESUMO deve estar localizada na margem
%% superior da folha e centralizada, e a referência, alinhada
%% à margem esquerda;
%% O alinhamento é justificado para o texto do resumo,
%% que inicia com parágrafo, e para as palavras-chave.

%% Espaçamento -> A palavra RESUMO deve ser separada da referência por
%% duas linhas em branco de 1,5;
%% Espaço 1 na referência e no resumo e, nas palavras-
%% chave, espaço 1,5.

%% Formato do papel,
%% orientação e margens -> Conforme especificado na seção 1.1.

%% Pontuação -> As palavras-chave devem ser separadas por ponto e
%% terminadas por ponto.

\begin{resumo}
\begin{SingleSpace}

\noindent
\begin{flushleft}
\entradaAutor{}. \textit{\imprimirtitulo}. \the\year. \pageref{LastPage} f. Trabalho de Conclusão de Curso (Graduação em Engenharia de Computação) - Instituto Politécnico, Universidade do Estado do Rio de Janeiro, Nova Friburgo, \the\year.
\end{flushleft}
\vspace{\onelineskip}

\setlength{\parindent}{1.3cm}

Os rompimentos das barragens de rejeitos em Mariana (2015) e Brumadinho
(2019) expuseram fragilidades críticas nos sistemas de monitoramento e
alerta do setor mineral brasileiro, resultando em centenas de mortes e
em danos ambientais de proporções históricas. Diante desse contexto,
este trabalho propõe o desenvolvimento de um sistema automatizado de
monitoramento e alerta para barragens de rejeitos de mineração, baseado
em tecnologia IoT, com o objetivo de mensurar continuamente a umidade do
solo, integrar dados meteorológicos externos e calcular um índice de
risco geotécnico em tempo real. O sistema é composto por um protótipo de
\textit{hardware} embarcado, desenvolvido sobre o microcontrolador
ESP8266 NodeMCU com sensor higrômetro resistivo, LEDs indicadores e
\textit{buzzer}, programado em C++ no ambiente Arduino IDE. O
\textit{firmware} integra as APIs meteorológicas BNDMET/DECEA e
OpenWeatherMap para obtenção de dados de precipitação histórica e
previsão pluviométrica, e transmite as leituras ao \textit{backend} em
intervalos adaptativos conforme o nível de risco calculado. O modelo de avaliação de risco integrado, desenvolvido pela autora,
fundamenta-se na abordagem multidimensional de \citeonline{rico2008mine},
segundo a qual 39\% dos 147 casos documentados de ruptura de barragens
de rejeitos resultaram da combinação de múltiplos fatores, e não de
uma causa isolada. Partindo dessa premissa, o modelo calcula um fator
de risco por meio da soma ponderada de sete variáveis que expressam
condições geotécnicas --- umidade do solo e taxa de variação da
umidade --- e meteorológicas --- precipitações acumuladas em três
janelas temporais, previsão pluviométrica e queda de pressão
atmosférica. O limiar de 24 horas apoia-se em \citeonline{rico2008mine}
e nos critérios do Sistema Alerta Rio \cite{sistemaalertario2025}, o de
30 dias foi definido a partir do histórico da estação D6594, e os
demais limiares e pesos foram arbitrados pela autora. O modelo inclui
ainda um mecanismo de amplificação, ativado quando a saturação elevada
do solo coincide com previsão de chuva moderada ou superior. A plataforma de \textit{software}
adota arquitetura cliente-servidor em três camadas: \textit{firmware}
embarcado, \textit{backend} desenvolvido em Node.js com TypeScript e
banco de dados PostgreSQL com extensão TimescaleDB para séries temporais,
e \textit{frontend} em React/Next.js com interface de monitoramento em
tempo real, histórico de leituras, gestão de usuários e central de
alertas. O sistema classifica o risco em três níveis (verde, amarelo e vermelho), 
sinalizados fisicamente pelos LEDs e pelo \textit{buzzer}
no dispositivo embarcado e exibidos na plataforma \textit{web}. O
administrador dispõe de uma central de alertas que permite o envio manual
de notificações por e-mail em massa aos usuários cadastrados, acompanhadas
do relatório do ciclo de análise correspondente. A solução proposta
demonstra que é possível desenvolver, com hardware de baixo custo e
tecnologias de código aberto, um sistema funcional de monitoramento
contínuo de risco geotécnico, a qual pode incluir a medição e o uso de
outras variáveis, como o nível de lençol freático e a poropressão do
solo (piezômetro) e a velocidade, a direção e a profundidade de
deslizamentos (inclinômetro), contribuindo para a prevenção de
acidentes em barragens de rejeitos de mineração.

\vspace{\onelineskip}
\noindent Palavras-chave: Barragens de rejeitos; ESP8266; Internet das Coisas; Monitoramento geotécnico; Sistemas de alerta.

\end{SingleSpace}
\end{resumo}